\talkslide{ipo}{DPO vs IPO loss ablation}{Timing: provisional}{IPO provides an alternative when preferences need not come from a single scalar Bradley--Terry reward.}{
\lead{Alternative formulation that doesn\textquotesingle t assume Bradley--Terry model}
\[
\mathcal L_{\rm IPO}=\mathbb E\left[
\left(\log\frac{\pi_\theta(y_w\mid x)\pi_{\rm ref}(y_l\mid x)}
{\pi_{\rm ref}(y_w\mid x)\pi_\theta(y_l\mid x)}
-\frac{1}{2\beta}\right)^2\right]
\]

\begin{statements}\setlength{\itemsep}{.3cm}
\item The formulation can represent \textbf{non-transitive preferences}
\item The squared loss targets a \textbf{finite likelihood-ratio margin}
\item IPO-MD supplies the motivation for \textbf{smoothed sampling}
\end{statements}
\paperrefs[14.7cm]{ipo,ipomd}
}
\speak{IPO provides an alternative when preferences need not come from a single scalar Bradley--Terry reward. The displayed generic online IPO loss targets a finite log-likelihood-ratio margin, rather than pushing it indefinitely. The game version couples attacker and defender preference problems and draws on IPO-MD. This addresses an over-optimization concern in the formulation; it is not a universal guarantee against overfitting or judge errors. The equation uses the background online IPO convention from Appendix B.2. Source detail: Azar et al. (2024); Calandriello et al. (2024); AdvGame, Appendix B.2, Eq. (34).}
